% 第 1 章 Introduction（原书 1-1 … 1-4）
\bichapter{引言}{Introduction}
\label{chap:intro}

\section{全芯片静态与动态电源完整性}
\subsection*{Full-chip Static and Dynamic Power Integrity}

ANSYS RedHawk\texttrademark\ 电源完整性（power integrity）解决方案是一款全芯片、基于单元（cell-based）的电源/地设计与验证产品，内建集成 SPICE，可从设计流程早期贯穿验证与签核（sign-off），覆盖静态与动态电源完整性。

静态 IR drop 依据电源分布的平均值确定压降，本质上只是解的直流（DC）分量。更关键的是交流（AC）效应：单元、时钟、存储器、IP 与 I/O buffer 等开关事件之间的时序关系，以及电容与电感对全芯片电源完整性的影响。

RedHawk 产品套件包含两个核心应用：
\begin{itemize}
  \item \textbf{RedHawk-S}（静态）：支持全芯片静态功耗、IR drop 与 EM（electromigration，电迁移）分析。
  \item \textbf{RedHawk-EV}（静态与动态）：支持全芯片 Vectorless Dynamic\texttrademark\ 功耗与电压分析，包含片上电感、封装 RLC 模型，以及去耦电容（decap）分析与优化。
\end{itemize}

此外，RedHawk 支持对先进低功耗设计的精确分析，例如 MTCMOS（电源门控）、时钟门控（clock gating）、多路 Vdd/Vss，以及衬底偏置（substrate biasing）。文中“RedHawk”可互换指代 RedHawk-S 或 RedHawk-EV，取决于所获许可的产品捆绑。

RedHawk 提供的综合物理电源完整性方案涵盖多项突破性技术，包括在基于单元的电源方案中达到晶体管级分析精度。其底层技术之一是 \textbf{Vectorless Dynamic} 统计分析引擎：无需功能向量即可在全芯片上计算现实的最坏开关场景。该技术由静态时序分析（STA）时序窗（timing window）以及经 SPICE 表征的 Apache Power Library（APL）电流波形驱动，可为电源网上每个实例生成准确的电压波形。峰值压降严重的区域被识别为热点（hotspot），可通过新增 decap 布局、网格尺寸调整、补充电源布线，和/或在电源热点处替换关键单元实例等方式修复。在识别电源网络热点后，RedHawk 可帮助判断如何修改网格、调整 decap 布局，或在所需位置替换单元实例，以降低压降。

RedHawk 与工业标准的工艺、设计与库描述文件格式完全互操作。为支持复杂 SoC 设计，RedHawk 接受混合输入：标准单元可用 LEF/DEF，存储器、flip-chip bump、I/O、cover macro 与 IP block 可用 GDSII。此外，RedHawk 支持以 APL 中的 SPICE 表征波形数据增强 Synopsys \texttt{.lib} 中的静态数据，从而保证“真正的”动态精度。

精度是 RedHawk 物理电源流程的主要目标之一；同等重要的是为设计者提供全芯片容量与易用性。RedHawk 在单一产品中实现：全芯片 RLC 电源/地提取、瞬态仿真、电迁移分析、静态与动态功耗及压降分析、压降对时钟抖动与延时等关键时序参数的影响，以及电源完整性修复与优化。设计者可据此在整个设计流程中快速评估物理实现决策对 SoC 功耗与时序的影响，并确保电源性能与签核准确。

\section{在设计流程中使用 RedHawk}
\subsection*{Using RedHawk in the Design Flow}

RedHawk 与工业标准格式完全兼容，可方便地嵌入现有 ASIC 厂商与 COT 流程。如图~\ref{fig:rh-design-flow} 所示，其物理电源方法学可集成到芯片设计的三个主要阶段：设计规划（Design Planning）、设计开发（Design Development）与设计验证（Design Verification）。以下各节说明 RedHawk 在这些设计阶段中的用法。

RedHawk 支持两种动态分析方法学，取决于所处流程阶段：
\begin{enumerate}
  \item 基于 \texttt{.lib} 文件数据：可在早期获得动态热点反馈。
  \item 基于 Apache Power Library（APL）的动态分析：在验证阶段基于表征后的单元电流波形，提供晶体管级精度。
\end{enumerate}

\subsection{设计规划}
\subsubsection*{Design Planning}

在此阶段，RedHawk 支持静态 IR drop 的早期分析，以及基于 \texttt{.lib} 的动态热点估计。此时可用线负载模型（如 net-to-gate 电容比）计算功耗分布。通常，IR drop 与 EM 的图形化结果可揭示电源/地网格结构的薄弱点或电源/地 pad 不足。电源网格分配问题可通过线宽调整，或增加 strap/via 等方式缓解；此时改用 flip-chip 封装也可能改善静态 IR 电源分配。此外，通过对尚不可得的存储器及其他内容做 black-box 处理，可在单个 block 或整个设计上标出因同时开关导致的动态热点。

\subsection{设计开发}
\subsubsection*{Design Development}

详细单元布局完成后，设计常以混合格式表示：逻辑核用 LEF/DEF，存储器与 flip-chip 电源 bump 用 GDSII。若信号网的 SPEF/DSPF 尚不可用，RedHawk 使用 Steiner tree 线长估计计算功耗分布。来自 STA 的 slew 信息可进一步提高功耗分析精度。可识别潜在 EM 问题。一旦识别全芯片动态热点区域，可快速进行 what-if 分析，评估封装模型电感与 intentional decap 的影响。尽早分析 decap，可确保在布局阶段而非布线之后分配合适的面积与位置。Decap 应尽量靠近受害点（victim），以有效抑制电流尖峰。

\figplaceholder{Figure 1-1}{在设计全过程中使用 RedHawk}
\label{fig:rh-design-flow}

图~\ref{fig:rh-design-flow} 概括了各设计阶段与 RedHawk 的主要活动：

\begin{itemize}
  \item \textbf{设计规划}（初始布局与电源网格设计）：
  \begin{itemize}
    \item 评估静态 IR 分布与 EM 问题
    \item 估计动态热点
    \item 评估并修复电源网格设计与电流需求不平衡
    \item 使用 \texttt{*.lib} 进行功耗分析
  \end{itemize}
  \item \textbf{设计开发}（详细布局与优化）：
  \begin{itemize}
    \item 执行基于 LIB 的动态热点分析与修复
    \item 探索并修改去耦电容
    \item 评估并修复封装与时钟电源问题
  \end{itemize}
  \item \textbf{设计验证}（详细布线、提取、LVS/DRC）：
  \begin{itemize}
    \item 执行基于 APL 的精确动态热点分析与修复
    \item 优化 decap 使用与防护
    \item 分析动态电压及其对时序的影响
    \item 评估并修复剩余电源相关问题
    \item 以接近 SPICE 签核精度验证设计
  \end{itemize}
\end{itemize}

\subsection{设计验证}
\subsubsection*{Design Verification}

在设计验证阶段，RedHawk 使用提取得到的 DSPF/SPEF 中信号网的详细阻容负载，以及 STA 给出的最终 slew 与 delay。验证阶段的关键数据是每个单元库实例的瞬态电流波形；这些波形表征于 Apache Power Libraries（APL）中，使单元级验证达到晶体管级精度。最终验证应使用 APL，以充分发挥 RedHawk 的晶体管级精度，并确保 decap 被最优放置以保护动态热点。流片前还需对动态电压引起的时钟 skew 与时序延时进行最终验证。

\section{小结}
\subsection*{Summary}

由于功耗密度升高、供电电压降低以及频率升高，全芯片动态与静态电源完整性已成为 65\,nm 及以下工艺节点设计的关键挑战之一。动态功耗与峰值压降属于瞬态现象，难以分析与修正，其对芯片时序与良率的影响日益受到关注。

主要 ASIC/IC 半导体厂商与 COT 公司可借助 ANSYS RedHawk 工具设计并验证超大规模设计的动态电源完整性。在 65\,nm 及更先进节点，动态电源完整性已成为设计签核要求。

一套完整的、基于单元的全芯片电源完整性分析方法学必须覆盖：
\begin{itemize}
  \item Vectorless Dynamic 压降分析与验证（含对时序与时钟树的影响）
  \item 去耦电容的分析与优化
  \item 面向单元与宏的波形级动态电源库
  \item 电源网格薄弱点分析、优化与修复
  \item 片上与片外电感噪声评估
  \item 每实例多路 Vdd/Vss
\end{itemize}

RedHawk 的物理电源方案在更快、高容量且易用的设计环境中提供上述全部能力。设计团队可在整个流程中放心部署其全芯片物理电源方法学，按计划流片，并以有保障的电源完整性交付硅片。
